\documentclass[12pt]{article}

\usepackage[margin=1in]{geometry}
\usepackage{setspace}
\usepackage{cite}
\usepackage{hyperref}

\title{Analysis of Tumor-Normal Single-Cell Ecosystems Across Cancer Types}
\author{Your Name \\ CISC 5790 Data Mining}
\date{}

\begin{document}

\maketitle

\onehalfspacing

\section{Introduction}

Understanding tumor biology at the single-cell level has become a central focus in cancer research. 
Traditional bulk sequencing methods average signals across many cells, which can obscure important cellular diversity.

In this study, the authors perform a large-scale single-cell analysis across nearly one thousand tumors spanning 30 cancer types 
\cite{kang2024systematic}. Their goal is to better understand tumor microenvironments and how cancer cells interact with surrounding normal and immune cells.

\section{Biological Context}

Tumors are not homogeneous masses of identical cells. Instead, they consist of a complex ecosystem of:
\begin{itemize}
    \item Cancer cells with diverse genetic profiles
    \item Immune cells responding to the tumor
    \item Stromal and supporting tissue cells
\end{itemize}

Single-cell technologies allow researchers to resolve these populations individually, providing a more detailed view of tumor biology.

\section{Biological Question}

The central question of this study is:

\textit{How do tumor and normal cells interact across different cancer types at single-cell resolution, and what patterns emerge across a large cohort of tumors?}

The authors aim to:
\begin{itemize}
    \item Characterize cellular composition of tumors
    \item Compare tumor ecosystems across cancer types
    \item Identify conserved and cancer-specific patterns
\end{itemize}

\section{Methods}

The study uses a large-scale single-cell RNA sequencing dataset and applies computational analysis to dissect tumor ecosystems.

Key steps include:
\begin{itemize}
    \item \textbf{Data collection:} Single-cell RNA-seq data from ~1000 tumors
    \item \textbf{Preprocessing:} Quality control, filtering, normalization
    \item \textbf{Clustering:} Grouping cells by gene expression profiles
    \item \textbf{Annotation:} Assigning biological identities to clusters
    \item \textbf{Cross-cancer comparison:} Identifying shared and unique patterns
\end{itemize}

\section{Results}

The authors identify several key findings:
\begin{itemize}
    \item Tumors exhibit highly diverse cellular ecosystems
    \item Immune infiltration patterns vary significantly across cancer types
    \item Certain cell states are conserved across multiple cancers
    \item Tumor microenvironment composition is linked to disease behavior
\end{itemize}

These findings suggest that tumor biology cannot be fully understood without considering single-cell heterogeneity.

\section{Discussion}

This study demonstrates the power of large-scale single-cell analysis in uncovering biological complexity.

However, challenges include:
\begin{itemize}
    \item High computational cost of single-cell analysis
    \item Data integration across multiple sources
    \item Biological interpretation of high-dimensional clusters
\end{itemize}

Despite these challenges, the study provides a comprehensive atlas of tumor ecosystems.

\section{Conclusion}

This work significantly advances our understanding of tumor heterogeneity by analyzing nearly 1000 tumors across 30 cancer types at single-cell resolution.

It highlights the importance of computational biology in revealing patterns that are not visible through traditional bulk sequencing approaches.

\bibliographystyle{plain}
\bibliography{references}

\end{document}
